\documentclass{article}

% NeurIPS 2026 main track, anonymized double-blind submission.
% Switch to \usepackage[preprint]{neurips_2026} for an arXiv-style preprint
% with author names visible. Use [main, final]{neurips_2026} only after
% acceptance for the camera-ready.
\usepackage{neurips_2026}

\usepackage[utf8]{inputenc}
\usepackage[T1]{fontenc}
\usepackage{hyperref}
\usepackage{url}
\usepackage{booktabs}
\usepackage{amsfonts}
\usepackage{amssymb}
\usepackage{amsmath}
\usepackage{nicefrac}
\usepackage{microtype}
\usepackage{xcolor}
\usepackage{graphicx}

% Compact symbol shortcuts; the source uses Unicode operators inline
% (\teArrow, \teImplies, etc.) which pdflatex needs as macros.
\newcommand{\teArrow}{\ensuremath{\rightarrow}}
\newcommand{\teImplies}{\ensuremath{\Rightarrow}}
\newcommand{\teTherefore}{\ensuremath{\therefore}}
\newcommand{\teBecause}{\ensuremath{\because}}
\newcommand{\teAnd}{\ensuremath{\wedge}}
\newcommand{\teOr}{\ensuremath{\vee}}
\newcommand{\teApprox}{\ensuremath{\approx}}
\newcommand{\teNeq}{\ensuremath{\neq}}

\title{The Architecture of Errors:\\Logarithmic Mode Discovery and Polylogarithmic Intervention Budgets for Long-Context LLM Reliability}

\author{%
  Anonymous Author(s)\\
  Affiliation\\
  Address\\
  \texttt{email}\\
}

\begin{document}

\maketitle

\begin{abstract}
Prior work on length-decay benchmarks~\citep{arbuzov2025beyond} argued that LLM reliability does not decay exponentially with output length because errors concentrate at sparse ``key tokens'' rather than spreading uniformly. This paper extends that framework along an orthogonal axis: LLM errors are not only sparse, they are \emph{repetitive}. The same handful of failure patterns recur across models, datasets, and domains. Inside the small set of key tokens, a latent fraction $\beta$ produce \emph{hard} failures; those hard failures cluster into a finite catalogue of recurring modes whose count grows slowly with the number of observed failures. Under an empirical postulate of logarithmic mode discovery---calibrated against endpoint counts from three published taxonomies (ErrorAtlas, MWPES, HumanEval-error-categorisation), with a direct subsample-discovery curve still outstanding---the intervention budget required to bound \emph{per-hard-token} residual error scales polylogarithmically in sequence length. Sequence-level reliability targets are strictly tighter and approach full-catalogue coverage as the number of hard decisions grows---an asymmetry we surface explicitly rather than burying. Twelve intervention categories from prior work, plus 28 GOLD-tier capability-elimination citations across six axes (Appendix~\ref{app:taxonomy}), are consistent with a small library on the order of tens of interventions covering the head of the failure distribution in the per-hard-token regime. We analyze the polylog conclusion symbolically under several cluster-count laws (logarithmic, Heaps, saturating; Appendix~\ref{app:heaps}) and find it robust across the family, though the specific doubly-logarithmic rate is the optimistic special case. Prominent steep-decay counter-evidence~\citep{dziri2023faith,kuratov2024babilong,kwa2025horizon}, re-audited in Appendix~\ref{app:reaudits}, decays over variables other than raw token length---compositional graph size, fact count, log-time horizon. This relocates rather than dissolves the practical worry: the framework identifies which interventions help and which do not, but does not claim the regimes where $k_{\text{hard}}$ grows fast are thereby made easy.
\end{abstract}

\section{Introduction}
\label{sec:intro}

The standard worry about long-context autoregressive generation is exponential. If every token has independent error probability $e$, then after $n$ tokens the chance of a fully correct output is $(1-e)^n$, which collapses to zero for any nontrivial $e$ and large enough $n$~\citep{lecun2023doomed,dziri2023faith}. Earlier work in this series~\citep{arbuzov2025beyond} argued that this worry is misplaced because errors are not distributed uniformly across tokens: only $5$--$10\%$ of tokens are ``key''---genuinely dependent on long-range context~\citep{fang2025perplexity}---while the remaining majority become nearly deterministic once enough context accumulates. Replacing $(1-e)^n$ with the two-rate model $P(\text{correct}) = (1-e_{\text{key}})^{k} (1-e_{\text{non}})^{n-k}$ with $k \ll n$ scaling sublinearly in $n$ recovers the observed long-context coherence and converts the reliability question from ``how does $n$ grow?'' to ``how does $k$ grow?''.

\paragraph{From sparse tokens to recurring patterns.}
This paper takes the next step. Sparsity tells us \emph{where} errors live; the natural follow-up is \emph{what} they are. Recent failure-mode atlases supply a striking empirical answer: errors are not only sparse, they are also \emph{repetitive}. ErrorAtlas~\citep{ashurytahan2026erroratlas}, covering 83 models across 35 datasets, organises observed failures into 17 named categories sorted by prevalence, with a long-tailed distribution heavily concentrated in the head. In code, two error types (\texttt{AssertionError} and \texttt{NameError}) cover 86.35\% of failures on HumanEval across 14 LLMs, replicated across 23 models on HumanEval Pro and MBPP Pro~\citep{wen2024fixing}. In math, MWPES-300K~\citep{sun2025mwpes} categorises 304,865 errors from 15 LLMs across four math word-problem datasets and reports that dataset characteristics shape error patterns systematically. The same picture appears in multi-hop QA~\citep{zhang2026weakestlink}, agentic tool use~\citep{cemri2025mast}, and retrieval-augmented generation~\citep{wood2024acurai}. These observations suggest a third layer in the reliability architecture, complementing key-token sparsity (Layer 1) and within-key stratified-manifold structure (Layer 2 in~\citealt{arbuzov2025beyond}). Within the small set of key tokens, only a fraction $\beta$ produce \emph{hard} failures, and those hard failures cluster into a finite catalogue of recurring modes whose size grows much more slowly than the number of observed events.

\paragraph{Contributions.}
(i)~A three-layer formal framework---sparsity ($\alpha$), stratification ($\beta$), and bounded mode catalogue ($|C|$)---that partitions key tokens into easy and hard subsets and adds a recurring-pattern catalogue (\S\ref{sec:framework}). (ii)~An empirically calibrated postulate of logarithmic mode discovery; anchored against ErrorAtlas, HumanEval, and MWPES, the mode-rate parameter falls in the range $\sigma \in [0.87, 1.85]$, with $\sigma \approx 1.85$ adopted from the ErrorAtlas anchor as a conservative target. (iii)~A polylog bound on the per-hard-token intervention library size, $m \geq \lceil |C|^{1 - \varepsilon / e_{\text{hard}}} \rceil$, with the doubly-logarithmic rate as the optimistic special case under the postulate; the \emph{sequence}-level analogue is strictly tighter and approaches full-catalogue coverage as $k$ grows. (iv)~Evidence synthesis for clustering, cluster-selective interventions, and sublinear length scaling (\S\ref{sec:evidence}), drawing on $\sim$60 prior published results including a six-axis capability-elimination harvest with 28 GOLD-tier citations (Appendix~\ref{app:taxonomy}). (v)~A re-audit of the most-cited steep-decay results (Appendix~\ref{app:reaudits}) showing each decays over a variable other than raw token length.

\section{Related Work}
\label{sec:related}

\paragraph{Failure-mode taxonomies.}
ErrorAtlas~\citep{ashurytahan2026erroratlas}, MWPES-300K~\citep{sun2025mwpes}, the HumanEval categorisation of~\citet{wen2024fixing}, and the RFMDataset~\citep{guo2025rfmdataset} together establish that, at any given corpus size, the number of distinct named failure modes is small (typically 8--20) and the cumulative coverage of the top modes is high. Domain-specific taxonomies for multi-hop QA~\citep{zhang2026weakestlink}, multi-agent systems~\citep{cemri2025mast}, and tool-augmented agents~\citep{yao2024taubench} report the same pattern.

\paragraph{Targeted interventions.}
Each named cluster has been the subject of focused intervention research. Python execution closes most of the arithmetic-error cluster~\citep{gao2022pal,chen2022pot,gou2024tora}; constrained decoding eliminates format violations~\citep{suresh2025dingo,wang2025slot,dong2024xgrammar}; execution feedback closes most code-logic errors~\citep{shinn2023reflexion,huang2023agentcoder}; process reward models reduce step-level reasoning errors~\citep{wangp2024mathshepherd,lightman2024verify}; RAG strongly reduces hallucinations~\citep{wood2024acurai}; preference optimisation closes over-refusal~\citep{karaman2024porover}; structured uncertainty fixes most tool-call failures~\citep{suri2025clarification}. Two structural observations recur: each intervention is \emph{cluster-selective} (residuals belong to a different named cluster, not to the targeted one), and \emph{additivity is approximate but not perfect}~\citep{patel2026sixsigma,le2026interaction}.

\paragraph{Length-decay benchmarks.}
A parallel literature measures the \emph{shape} of the reliability decay curve. Mild-decay results---Loong~\citep{wangm2024loong}, GSM-$\infty$~\citep{zhou2025gsminfinite}, RULER~\citep{hsieh2024ruler}, anchor-based LLMs~\citep{pang2024anchor}---report log-linear, sigmoidal, or threshold-like decay inconsistent with smooth $(1-\varepsilon)^n$. Steep-decay results~\citep{dziri2023faith,kuratov2024babilong,kwa2025horizon,wan2026fano} are sometimes cited as exponential-compounding evidence. The apparent tension dissolves on careful reading: every steep-decay paper decays over a variable distinct from raw token length, which Appendix~\ref{app:reaudits} documents.

\paragraph{Gap.}
What is missing is a quantitative bridge between the small, repeating taxonomy catalogue and a polylog intervention budget. We provide that bridge in \S\ref{sec:framework}.

\section{Theoretical Framework}
\label{sec:framework}

\paragraph{Levels of $|C|$.}
LLM error analysis routinely conflates four distinct objects, which we keep separate. \emph{L1: failure events}---concrete observed errors. \emph{L2: empirical taxonomy categories}---researcher-chosen labels grouping L1 events; what published taxonomies report. \emph{L3: latent failure modes}---the unobserved underlying clusters L2 approximates. \emph{L4: capability axes / interventions}---the engineering unit (Python interpreter, constrained decoder, RAG). Throughout this section, $|C|$ refers to the L2 count, because L2 is what published taxonomies report; Postulate~1 is engineering-relevant because L4 is coarser than L2, and epistemically conditional because L2 is a noisy proxy for L3.

\subsection{$\beta$-stratification of key tokens}

The two-rate model of~\citet{arbuzov2025beyond} distinguishes $k$ key tokens (error rate $e_{\text{key}}$) from $n - k$ non-key tokens ($e_{\text{non}} \ll e_{\text{key}}$). Empirical atlases suggest that even within the key-token class errors are not uniformly distributed: most key tokens are ``decisions'' for which the model has stable representations; only a fraction concentrate the actual failures.

\textbf{Definition 1 (Hard fraction).} Partition the $k$ key tokens of a sequence into easy and hard subsets, $k_{\text{hard}} = \beta k$, $k_{\text{easy}} = (1-\beta) k$, with $\beta \in (0,1)$. Hard key tokens have an elevated error rate $e_{\text{hard}}$ corresponding to manifold-transition decisions in the sense of~\citet[\S3.2]{arbuzov2025beyond}; easy key tokens have $e_{\text{easy}} \teApprox e_{\text{non}}$. The composed reliability becomes
\[
P(\text{correct}) = (1 - e_{\text{hard}})^{\beta k} (1 - e_{\text{easy}})^{(1 - \beta) k} (1 - e_{\text{non}})^{n - k}.
\]
We do not assume iid token errors: the three rates summarise \emph{conditional per-decision failure hazards} after conditioning on prior decisions being correct, and grouping by stratum yields the multiplicative survival expression. Because $e_{\text{easy}} \teApprox e_{\text{non}}$ once context accumulates, the failure rate is dominated by the $\beta k$ hard tokens, and we treat $e_{\text{hard}}$ as the load-bearing quantity. The parameter $\beta$ is \emph{latent}: failure-mode atlases report $\Pr(\text{category} \mid \text{error occurred})$, which does not identify $\Pr(\text{hard} \mid \text{key-token decision})$. We carry $\beta$ symbolically.

\subsection{Empirical postulate: logarithmic mode discovery}
\label{sec:postulate}

The $\beta k$ hard-token failures cluster into a finite set of recurring patterns indexed $c = 1, \ldots, |C|$. The natural question is how $|C|$ grows with the number of observed failures. Two candidates exist in the type--token literature: Heaps' law (power-law, $|C| \approx K \cdot k_{\text{hard}}^{b}$ with $b \in [0.4, 0.6]$~\citep{manning2008ir}), and logarithmic discovery. \textbf{Zipfian rank-frequency does not imply logarithmic growth}: Heaps' law is the correct type--token consequence of Zipf-distributed events, and it is power-law, not logarithmic. We therefore state logarithmic mode discovery as an empirical postulate defensible by direct measurement, not as a theorem.

Two sample-size variables matter and conflating them is the seam that lets reviewers in. Let $h(n) = \beta k(n)$ denote the number of hard-token decisions in a single sequence of length $n$; let $T$ denote observed hard-failure events in a corpus used to discover the domain catalogue. We write $C_D$ for the full reachable catalogue inside domain patch $D$, $C_{\text{seen},D}(T)$ for the subset discovered after $T$ failures, and $C_{\text{active},D}(n)$ for the subset a single sequence of length $n$ can activate.

\textbf{Postulate 1 (Patch-indexed catalogue discovery).}\label{post:logmode} \emph{Within a fixed application domain $D$, the number of named recurring failure modes discovered after $T$ sampled hard-failure events is bounded by}
\[
|C_{\text{seen},D}(T)| \leq \min\bigl(A_D + \sigma_D \ln T,\ |C_D|\bigr), \qquad \sigma_D > 0,\ A_D \geq 0.
\]
\emph{The cap $|C_D|$ is the domain-imposed ceiling. The postulate concerns catalogue discovery, not the number of hard decisions inside a single sequence.}

\textbf{Assumption 2 (Per-sequence active-mode exposure).} \emph{For a single sequence of length $n$ in domain $D$,}
\[
|C_{\text{active},D}(n)| \leq \min\bigl(A'_D + \sigma'_D \ln h(n),\ |C_D|\bigr).
\]
\emph{Primed constants are distinct from those of Postulate 1; corpus discovery and per-sequence activation need not share the same rate. Proposition~1's sequence-length scaling uses $C_{\text{active},D}(n)$; domain-library budgeting uses $C_D$.}

\textbf{Empirical defense.} ErrorAtlas: $k_{\text{hard}} \gtrsim 10^4$, $|C| = 17$, giving $\sigma \teApprox 17 / \ln(10^4) \teApprox 1.85$ with $A=0$. HumanEval~\citep{wen2024fixing}: $|C| \in [8, 12]$, $\sigma \in [0.87, 1.30]$. MWPES: $|C| \in [15, 20]$, $\sigma \in [1.2, 1.6]$. Across the three anchors, $\sigma \in [0.87, 1.85]$; we carry $\sigma \teApprox 1.85$ as the conservative target. Endpoint counts at a single corpus scale are not discovery curves; the subsample-discovery measurement remains an explicit empirical test we flag (\S\ref{sec:limitations}, Appendix~\ref{app:heaps}).

\textbf{Capabilities are coarser than error classes.} Postulate~1's $\sigma \teApprox 1.85$ is calibrated against \emph{named error categories}, but a single capability axis often targets several L2 categories at once. A Python interpreter removes the execution-error component of arithmetic, unit conversion, simple counting, list manipulation, and date arithmetic; a constrained decoder eliminates by construction the format-violation cluster and the structural component of ``missing required element.'' Postulate~1 is therefore a conservative input for the engineering question; the formal coverage problem is weighted set cover, with the ranked-category form of Proposition~1 as a tractable proxy. An alternative power-law (Heaps) formulation and symbolic-form sensitivity analysis across candidate cluster-count laws are given in Appendix~\ref{app:heaps}.

\textbf{Domain patches cap the catalogue.} A model deployed in a fixed application domain occupies a bounded sub-region of the stratified manifold---a \emph{domain patch} $P_D$~\citep{arbuzov2025beyond,li2025stratified}. Empirical support: embedding-space intrinsic dimensions of $\sim 5$--$15$ across domains, well below ambient $\sim 10^3$~\citep{park2024linear,li2025stratified}; cross-domain accuracy spreads of $20$--$70\%$ on a single base model~\citep{wangy2024mmlupro}; log-popularity thresholds below which scaling fails~\citep{mallen2023popqa,kandpal2023longtail}. None of these \emph{measure} $|C_D|$ directly; together they motivate the patch-indexed hypothesis. The patch makes logarithmic or saturating discovery dynamics plausible inside fixed domains; $\sigma_D$, $A_D$, $\beta_D$, $|C_D|$ are all domain-indexed.

\subsection{Coverage by a targeted intervention library}
\label{sec:coverage}

Given a catalogue of $|C|$ recurring modes, an intervention library of size $m$ covers the top-$m$, leaving $|C|-m$ uncovered. We adopt the log-coverage form
\[
F(m; |C|) = \min\!\left(1,\ \frac{\ln m}{\ln |C|}\right)
\]
as a \emph{postulated} closed expression respecting the $\min(1,\cdot)$ cap. For our ErrorAtlas anchor ($|C|=17, m=5$), $F(5;17) \teApprox 56.8\%$ and $F(10;17) \teApprox 81.3\%$. This is an empirical best-fit choice, not derived from Zipf: ErrorAtlas sorts its 17 categories by prevalence but does not publish a single cross-domain cumulative top-$m$ curve at this writing; alternative cumulatives (Zipf--Mandelbrot, truncated power law, saturating exponential) match the same qualitative shape with comparable freedom. Appendix~\ref{app:heaps} shows the polylog conclusion survives across them. Two caveats: (a)~$F$ counts named L2 categories covered, not L4 capability axes deployed, so real-deployment residual error is often lower at any given $m$; (b)~we adopt $F(0;|C|) := 0$ and $F(1;|C|) := p_1$ (mass of the most prevalent mode) since the log form is not anchored at $m=1$. After deploying the top-$m$ library, the residual per-hard-token error rate is $e_{\text{res}}(m) = (1 - F(m; |C|))\, e_{\text{hard}}$.

\subsection{Main proposition: conditional polylogarithmic intervention budget}
\label{sec:proposition}

We compose the three ingredients with the two-rate model. For sequence-length scaling we use Assumption~2, not Postulate~1: the empirical sparsity $\alpha = k/n \in [0.05, 0.10]$ is a \emph{finite-length} statement, and under~\citet[\S3.1]{arbuzov2025beyond}'s $k \sim \log n$ regime, $\alpha$ decays as $O(\log n / n)$.

\textbf{Conditionality.} The bound below is engineering math, not a theorem about LLMs. It holds conditional on (i)~the log-coverage form of \S\ref{sec:coverage} on its declared domain $m \geq 2$, $|C_{\text{eff}}| \geq 2$; (ii)~the residual target $\varepsilon \in (0, e_{\text{hard}})$; (iii)~Assumption~2 for any sequence-length rate claim. A different cumulative fitted to the same anchors preserves the qualitative polylog conclusion (Appendix~\ref{app:heaps}) but changes the exponent.

\textbf{Proposition 1 (Per-Hard-Token Intervention Budget).}\label{prop:budget} \emph{Let $\varepsilon \in (0, e_{\text{hard}})$ be a target per-hard-token residual error rate for sequences in domain $D$, with $|C_{\text{eff}}| \geq 2$, where $C_{\text{eff}} = C_{\text{active},D}(n)$ for per-sequence scaling or $C_{\text{eff}} = C_D$ for full domain-library deployment. Under the coverage form of \S\ref{sec:coverage}, the smallest intervention library satisfying $e_{\text{res}}(m) \leq \varepsilon$ has size}
\[
m \geq \left\lceil |C_{\text{eff}}|^{1 - \varepsilon / e_{\text{hard}}} \right\rceil.
\]
\emph{If $C_{\text{eff}} = C_{\text{active},D}(n)$ and $k(n) = \Theta(\log n)$, then in the pre-cap regime $m = O\bigl((\log\log n)^{1 - \varepsilon / e_{\text{hard}}}\bigr)$, i.e., the intervention budget grows doubly logarithmically with sequence length. In the cap regime, or for full domain-library deployment, $m \geq \lceil |C_D|^{1 - \varepsilon / e_{\text{hard}}} \rceil$, independent of $n$.}

\emph{Proof.} From \S\ref{sec:coverage}, $e_{\text{res}}(m) = (1 - F(m; |C_{\text{eff}}|))\,e_{\text{hard}}$. Requiring $e_{\text{res}}(m) \leq \varepsilon$ gives $F(m; |C_{\text{eff}}|) \geq 1 - \varepsilon / e_{\text{hard}}$. Since $\varepsilon < e_{\text{hard}}$, the required coverage lies below one, so the cap in $F$ is inactive; with $F = \ln m / \ln |C_{\text{eff}}|$, we obtain $m \geq |C_{\text{eff}}|^{1 - \varepsilon / e_{\text{hard}}}$, and taking the ceiling gives the stated form. Assumption~2 gives $|C_{\text{active},D}(n)| = O(\log\log n)$ when $k(n) = \Theta(\log n)$, before the cap. After the cap, $|C_{\text{eff}}| = |C_D|$ and the bound is independent of $n$. $\square$

We label this a \emph{Proposition} rather than a \emph{Theorem} to keep its conditional nature visible. The doubly-logarithmic rate is the \emph{optimistic special case}: per-sequence, pre-cap, under logarithmic mode discovery. Weaker (but still polylog) rates obtain under Heaps (Appendix~\ref{app:heaps}); cap regime is a domain-constant.

\textbf{Sequence-level versus per-hard-token bounds.} Proposition~1 bounds the \emph{per-hard-token} residual error. Sequence-level failure probability is strictly stronger. With non-hard-token survival factor $S_{\text{base}} = (1 - e_{\text{easy}})^{(1 - \beta) k}(1 - e_{\text{non}})^{n - k}$, requiring $(1 - e_{\text{res}})^{\beta k}\, S_{\text{base}} \geq 1 - \varepsilon_{\text{seq}}$ yields three regimes: (i)~if $S_{\text{base}} < 1 - \varepsilon_{\text{seq}}$, the binding constraint is non-key noise and hard-token interventions alone cannot meet the target---this is the back-door route by which the exponential-in-$n$ concern re-enters; (ii)~if $S_{\text{base}} \geq 1 - \varepsilon_{\text{seq}}$, the required library size is $m \geq \lceil |C_{\text{eff}}|^{1 - \tau_{\text{seq}} / e_{\text{hard}}} \rceil$ with $\tau_{\text{seq}} = 1 - ((1-\varepsilon_{\text{seq}})/S_{\text{base}})^{1/(\beta k)}$, and as the target tightens toward feasibility, $\tau_{\text{seq}} \to 0$ and $m \to |C_{\text{eff}}|$---essentially every named mode must be covered; (iii)~if $\tau_{\text{seq}} \geq e_{\text{hard}}$, the target is met without intervention. The widely-quoted ``$\teApprox 50$ patterns $\teArrow \teApprox 90\%$ reduction'' claim is a per-hard-token statement; the sequence-level analogue requires correspondingly more interventions. Production systems with one-shot reliability targets should plan for the stricter regime---and check whether they are in Regime (i), where catalogue-budgeting cannot help at all.

\section{Empirical Evidence}
\label{sec:evidence}

We collect evidence for three load-bearing claims: errors cluster into a small recurring set (A); each cluster is addressable by one targeted intervention (B); reliability decays sublinearly with output length (C).

\paragraph{Claim A: Failure-mode clustering.}
Across-domain concentration is strongest in ErrorAtlas~\citep{ashurytahan2026erroratlas}: 83 models $\times$ 35 datasets, $\teApprox 7{,}000$ failures, 17 named categories with a long-tailed, head-concentrated prevalence ordering. Domain-specific Paretos reproduce the shape: math (MWPES, top-4 categories dominant per model with strong cross-model overlap~\citep{sun2025mwpes}); code (\texttt{AssertionError}+\texttt{NameError} = 86.35\% on HumanEval, replicated across 23 models~\citep{wen2024fixing}); multi-hop QA (single dominant ``missing-evidence'' failure mode~\citep{zhang2026weakestlink}); agentic tool use (under 20 recurring modes~\citep{cemri2025mast,yao2024taubench}); RAG (handful of distinguishable hallucination types~\citep{wood2024acurai}). Cross-model invariance: ErrorAtlas reports a stable category hierarchy across 83 models; RFMDataset~\citep{guo2025rfmdataset} reports ``strikingly similar failure mode distributions'' across 10 advanced reasoning models; EDIT~\citep{dai2025edit} finds a $\teApprox 4.7\%$ key-step fraction stable across models. Theoretical support: \citet{schaeffer2025monkeys} prove that observed aggregate power-law scaling requires the per-problem success-rate distribution to be heavy-tailed near $p=0$---structurally related to Postulate~1.

\paragraph{Claim B: Cluster-selective capability interventions.}
A dedicated harvest yields 28 GOLD-tier citations across six independent capability axes (arithmetic, code execution, format/structure, perception/grounding, knowledge/RAG, verification), each independently confirmed by 3--9 citations. The full table and structural patterns (Pattern A by-construction, B strong-empirical-with-class-shift, C moderate-with-shift) appear in Appendix~\ref{app:taxonomy}. Headline numbers: PAL~\citep{gao2022pal} lifts GSM-Hard from 20.1\% to 61.5\% with residuals shifting to problem-comprehension; constrained decoding eliminates format violations by construction~\citep{suresh2025dingo,wang2025slot,dong2024xgrammar}; Reflexion~\citep{shinn2023reflexion} and AgentCoder~\citep{huang2023agentcoder} push HumanEval pass@1 from 80\% to 96.3\% with residuals in spec-misinterpretation; Math-Shepherd~\citep{wangp2024mathshepherd} raises Mistral-7B MATH from 28.6\% to 43.5\% with process supervision and verification reranking; Acurai~\citep{wood2024acurai} reports 100\% hallucination elimination on RAGTruth (a benchmark-conditional empirical result, not a general by-construction guarantee); POROver~\citep{karaman2024porover} lifts not-overrefused rate from 57.6\% to 82.1\%; SAGE-Agent~\citep{suri2025clarification} raises When2Call from 36.5\% to 65.2\%. The class-shift signature---post-intervention residuals belong to structurally different classes---is the empirical content of the cluster-orthogonality claim that underwrites Proposition~1's composition step. Capabilities are coarser than error classes: a single Python interpreter removes execution-error components from five named clusters; constrained decoding eliminates format and the structural component of ``missing required element'' jointly. A negative control~\citep{tian2024debugbench} confirms selectivity: execution feedback works for syntax/reference errors and is explicitly ``unhelpful for logic errors''---exactly where the framework predicts capability provisioning fails.

\paragraph{Claim C: Sublinear length scaling.}
Loong~\citep{wangm2024loong}: GPT-4o on Chain-of-Reasoning declines 81.6\% $\teArrow$ 32.9\% across $10$K $\teArrow$ $>$200K context bins---log-linear over $\log L$, decisively non-exponential per-token. GSM-$\infty$~\citep{zhou2025gsminfinite} fits a sigmoid with $R^2 > 0.98$; exponentially more inference compute yields only \emph{linear} AUC gains (direct signature of clustered, not independent, errors); DeepSeek-R1 maintains $10\%+$ accuracy at 130 reasoning operations, ruling out $(1-\varepsilon)^{130}$ for any reasonable $\varepsilon$. Compositionality gap~\citep{press2023compositionality}: $\teApprox 40\%$ gap between two-hop and product-of-one-hop accuracies, approximately \emph{constant} across the GPT-3 family---direct evidence against per-step independent compounding. Anchor LLMs~\citep{pang2024anchor}: $\teApprox 99\%$ K/V reduction with minor accuracy compromise---the effective number of distinct attention-key vectors needed is far smaller than $n$. Staircase patterns: Think-Prune-Train~\citep{costello2025think} raises Pass@1 while Pass@20 plateaus (the manifold-transition signature); RULER~\citep{hsieh2024ruler} reports threshold behaviour inconsistent with smooth $(1-\varepsilon)^n$. METR~\citep{kwa2025horizon} fits logistic-in-log-length, mathematically sublinear in length itself. Self-consistency~\citep{wangx2023selfconsistency} is consistent with---but does not prove---clustered structure: Condorcet aggregation of iid Bernoulli chains can yield similar magnitude gains, so we treat it as suggestive rather than diagnostic. Prominent steep-decay counter-evidence~\citep{dziri2023faith,kuratov2024babilong,wan2026fano,karpinska2024nocha} decays over variables distinct from raw token length; per-paper re-audits showing the relocation pattern are in Appendix~\ref{app:reaudits}.

\section{Practical Implications}
\label{sec:implications}

\paragraph{Reliability engineering is local patch coverage.}
Within a fixed domain patch, Proposition~1 makes reliability a small-catalogue engineering problem rather than an asymptotic scaling problem. A team should initially budget for a library on the order of tens of interventions, then refine from the local mode-discovery curve $C_{\text{seen},D}(T)$ and the empirical rank-frequency distribution. Available evidence is consistent with $\teApprox 50$ named interventions covering the head of the per-hard-token failure distribution in many measured domains, but this number is a planning prior, not a universal constant: the same base model in cardiology RAG, legal contract drafting, and code review yields three different libraries because $C_D$, $A_D$, $\sigma_D$ all change with the deployment domain.

\paragraph{Per-hard-token vs.\ sequence-level targets matter for SLA design.}
Per-hard-token residual error---the right metric for continuous-correction systems---has a polylog budget. Sequence-level failure probability---the right metric for one-shot systems---is strictly stricter and approaches full-catalogue coverage as $k$ grows. Production deployments should design SLAs to match the actual cost structure, not read the per-token result as the production SLA.

\paragraph{Library design as engineering, not asymptotic theory.}
An intervention library can be assembled module-by-module, in any order, as long as modules target distinct failure classes. Layer-separated interventions (constrained decoding + retrieval + process supervision + tool call) compose approximately additively in our data; the~\citet{le2026interaction} caveat applies only to interventions sharing a prompt channel. Library construction is therefore highly parallelisable.

\paragraph{Capability libraries are coarser than error-class libraries.}
Five named math classes (arithmetic, units, counting, list manipulation, date arithmetic) collapse under one Python interpreter; three code classes (SyntaxError, NameError, most TypeError) collapse under one execution-feedback loop; format violations and the structural component of ``missing required element''---together more than $20\%$ of ErrorAtlas---collapse under one constrained-decoder deployment. A team budgeting for $\teApprox 50$ interventions per domain can expect to need fewer than 50 capability-axis interventions while covering $\teApprox 50$ named clusters. The six axes of Appendix~\ref{app:taxonomy} are closer to the right unit of engineering accounting than the twelve clusters.

\section{Discussion}
\label{sec:discussion}

\paragraph{By-construction elimination as a boundary case.}
Seven of the 28 GOLD citations (Appendix~\ref{app:taxonomy}; constrained decoders~\citep{suresh2025dingo,dong2024xgrammar,dong2026xgrammar2,zhang2023tooldec,openai2024structured}, proof kernels~\citep{ren2025deepseekprover}, static syntax checks) achieve residual error rate equal to zero \emph{by mathematical construction}, not statistically. Constrained decoders set $P(\text{invalid token}) = 0$ at every generation step; the class of grammar-violating outputs is mathematically empty. In this regime, the polylog bound of Proposition~1 is loose: covered clusters contribute zero to residual error, not a small positive quantity. This is a strengthening result, not a counter-example---capability provisioning erases the error class entirely. Pattern A applies precisely to structural/verifiable classes (format, syntax, schema, invalid proofs), which are also the head of the heavy-tailed cluster frequency distribution. The strongest mechanism applies exactly where the catalogue is densest, compounding the polylog conclusion.

\paragraph{Counter-evidence relocates, not dissolves.}
Five prominent steep-decay papers~\citep{dziri2023faith,kuratov2024babilong,kwa2025horizon,wan2026fano,karpinska2024nocha} are routinely cited as exponential-compounding evidence. Every one decays over a variable distinct from raw token length: compositional graph size~\citep{dziri2023faith}, number of supporting facts~\citep{kuratov2024babilong}, log-time horizon~\citep{kwa2025horizon}, capacity threshold~\citep{wan2026fano}, evidence scope~\citep{karpinska2024nocha}. The apparent rapid decay is, in each case, in a quantity our framework already concentrates the action in ($k_{\text{hard}}$ and $|C|$), not in raw sequence length---a relocation, not a dissolution. These regimes remain hard; the framework's value is directing intervention toward capability provisioning along the actual decay axis rather than toward context-window expansion that does not help. Appendix~\ref{app:reaudits} walks through each paper.

\paragraph{Limitations.}
\label{sec:limitations}
Postulate~1 is empirical, not derived; a domain with genuinely Heaps-power-law mode discovery would invalidate the doubly-logarithmic special case while preserving the qualitative polylog conclusion. The coverage form of \S\ref{sec:coverage} is an empirical best-fit; readers should not interpret the numerical match against anchors as theoretical confirmation. Cluster-orthogonality is approximate; \citet{le2026interaction}'s caveat on prompt-channel interference tightens the bound in shared-channel settings. Domain narrowness: Postulate~1's empirical anchor rests on three taxonomies (ErrorAtlas, HumanEval, MWPES), all published 2025--2026 and covering general/code/math; the framework is untested on agentic workflows, long-running scientific reasoning, and multi-turn tool use over millions of tokens---precisely the regimes where $|C|$ may grow faster than logarithmically. The $\sigma \teApprox 1.85$ estimate is a single-point calibration, not a discovery-curve fit; a subsample-vs-distinct-modes plot has not been published for any LLM error taxonomy at this writing. $\beta$ is latent and we report no empirical range. Taxonomy granularity (the L2--L3 gap) is unmeasured. Patch-shift between deployments changes $A_D, \sigma_D, \beta_D, |C_D|$ simultaneously; libraries calibrated against one patch under-cover the next. Finally, the framework \emph{relocates} the difficulty of long-context reliability rather than resolving it: where $k_{\text{hard}}$ grows with task length (adversarial compositional structure, multi-hop chains, long agent horizons), reliability remains hard. The contribution is to name the on-axis intervention, not to make those regimes easy. A falsifiability test the framework passes: DebugBench~\citep{tian2024debugbench} explicitly catalogues which classes execution-feedback can and cannot repair, and the empirical split (works for syntax/reference; ``unhelpful for logic errors'') is exactly what the framework's selectivity claim predicts.

\section{Conclusion}
\label{sec:conclusion}

LLM reliability is often framed as an asymptotic scaling problem. Prior work in this line~\citep{arbuzov2025beyond} argued that the framing rests on a false uniformity assumption---errors concentrate at $\teApprox 5$--$10\%$ of tokens. This paper takes the next step: within that sparse set, errors are not only concentrated but also \emph{repetitive}, clustering into a finite catalogue whose size, under the empirical postulate of \S\ref{sec:postulate}, grows logarithmically---and under the conservative Heaps alternative (Appendix~\ref{app:heaps}), as a small power---of observed failures. The conditional consequence (Proposition~1) is that the intervention budget required to bound per-hard-token residual error scales polylogarithmically in sequence length within a fixed domain patch, and becomes a domain-constant once the patch ceiling $|C_D|$ is reached. Sequence-level reliability targets are strictly tighter. Available evidence is consistent with libraries on the order of tens of interventions covering the head of the per-hard-token failure distribution in many fixed domains. The exact number is patch-indexed and should be estimated from local discovery and rank-coverage curves, not assumed from cross-domain priors. This reframes the question from ``can we bound the growth of $n$-token error?'' to ``have we catalogued enough failure modes?'' The latter is finite and addressable. Direct empirical measurement of the mode-rate constant $\sigma$ on new domains---agentic, scientific, code-in-production---is the open empirical question this paper invites; future work flagged by~\citet[\S6]{arbuzov2025beyond} on attention-mechanism bounds and stratified-manifold geometry would replace the empirical postulate with a derived result.

\bibliographystyle{plainnat}
\bibliography{references}

\appendix

\section{Heaps Power-Law Variant and Cluster-Count Sensitivity}
\label{app:heaps}

A reader who prefers to derive cluster-count growth from standard Heaps' law rather than from Postulate~1 obtains a qualitatively similar result. Let $|C|(k_{\text{hard}}) \teApprox K \cdot k_{\text{hard}}^{b}$ with $b \in (0, 1)$. Canonical fits give $b \teApprox 0.5$ for natural-language vocabularies; failure-mode taxonomies plateau much more sharply (ErrorAtlas stabilises at $|C| = 17$ across $10^4{+}$ failures), implying a smaller effective $b \teApprox 0.2$ in our setting. Composing with $k = \Theta(\log n)$, we get $|C| = O((\log n)^{b})$, and Proposition~1 becomes
\[
m = O\!\left((\log n)^{b \cdot (1 - \varepsilon / e_{\text{hard}})}\right).
\]
This is still polylogarithmic in $n$ for any $b \in (0, 1)$ and any $\varepsilon < e_{\text{hard}}$. The paper's qualitative claim survives either choice of cluster-count law.

\paragraph{Symbolic-form sensitivity across candidate laws.}
The polylog conclusion depends on which cluster-count law one accepts; available evidence is consistent with multiple candidates because no subsample-discovery curve has been published for any LLM failure-mode taxonomy at this writing. We therefore report symbolic rates rather than fitted constants. With $h(n) = \beta k(n)$:

\begin{itemize}
\setlength{\itemsep}{1pt}
\item \textbf{Logarithmic:} $|C_{\text{active},D}(n)| = O(\log h(n))$. If $k(n) = \Theta(\log n)$, then $m = O\bigl((\log\log n)^{1 - \varepsilon / e_{\text{hard}}}\bigr)$.
\item \textbf{Heaps:} $|C_{\text{active},D}(n)| = O\bigl(h(n)^b\bigr)$ with $b \in (0, 1)$. If $k(n) = \Theta(\log n)$, then $m = O\bigl((\log n)^{b\,(1 - \varepsilon / e_{\text{hard}})}\bigr)$.
\item \textbf{Saturating:} $|C_{\text{active},D}(n)| \leq |C_D|$. Then $m = O\bigl(|C_D|^{1 - \varepsilon / e_{\text{hard}}}\bigr)$, constant in $n$ once the patch ceiling is reached.
\end{itemize}

The qualitative conclusion is robust \textbf{under the $k(n) = \Theta(\log n)$ regime}: $m$ grows more slowly than any positive power of $n$ under every candidate, and the directional claim (``a small library covers the head of the failure distribution in the per-hard-token regime'') survives. Only the exponent shifts: doubly-logarithmic under logarithmic discovery, $(\log n)^b$ with small $b$ under Heaps, constant in the cap regime. The doubly-logarithmic rate is the optimistic special case. If $k(n)$ grows as a positive power of $n$, the Heaps variant inherits that power and the polylog-in-$n$ language fails along that axis---the framework's intervention prescription still applies, but its asymptotic-rate framing does not.

Numerical constants require a measured discovery curve $C_{\text{seen},D}(T)$ or $C_{\text{active},D}(n)$. Existing taxonomies provide endpoint category counts at a single corpus scale (ErrorAtlas at $|C|=17$ for $\teApprox 10^4$ failures), not discovery curves. The subsample-discovery measurement remains the explicit empirical test that would either tighten the postulate or fall back to the Heaps variant. Until that measurement exists, the framework's headline rate should be read as \emph{polylogarithmic in the pre-cap regime, domain-constant in the cap regime}---with the specific exponent flagged as a falsifiability test rather than a fitted prediction.

\section{Full Failure-Mode Taxonomy and the 28 GOLD Citations}
\label{app:taxonomy}

The intervention literature provides at least one targeted countermeasure for each named cluster of \S\ref{sec:evidence}. Two organising granularities exist: at the capability level (six axes) the cluster-orthogonality property is most clearly visible; at the error-class level (twelve named clusters) the evidence aligns with the taxonomies of \S\ref{sec:evidence} but is finer-grained than the underlying capability mechanisms.

\paragraph{Six capability axes and three structural patterns.}
A dedicated harvest yields \textbf{28 GOLD-tier citations} across \textbf{six independent capability axes}: (i)~Arithmetic (Python/symbolic execution), (ii)~Code Execution (REPL/sandbox feedback), (iii)~Format/Structure (constrained decoding, FSMs, grammar engines), (iv)~Perception/Grounding (visual grounding for GUI, charts, tables), (v)~Knowledge/RAG (dense retrieval and citation grounding), (vi)~Verification (proof checkers, learned verifiers, classifier rerouting, process supervision). Each axis is independently confirmed by 3--9 GOLD citations. We stratify into three tiers by kind of evidence:

\textbf{Pattern A --- Hard guarantees (by-construction).} Seven citations achieve residual error rate equal to zero \emph{by construction}, restricted strictly to structural/verifiable classes: constrained decoders set $P(\text{invalid token}) = 0$ at every step~\citep{suresh2025dingo,zhang2023tooldec,dong2024xgrammar,dong2026xgrammar2,openai2024structured}; static syntax checks reject programs with a \texttt{SyntaxError} before execution~\citep{wen2024fixing}; proof kernels reject any output failing type-checking~\citep{ren2025deepseekprover}. The class of grammar-violating outputs is mathematically empty under these mechanisms.

\textbf{Pattern B --- Strong empirical reductions with class-shift signature.} Roughly fourteen citations report empirical reductions of $80$--$100\%$ in a named error class, with the post-intervention failure log dominated by structurally different residual classes. Program-of-Thoughts on GSM8K~\citep{chen2022pot}: calculation errors drop from $30\%$ of failures to $0\%$, residuals are $62\%$ reasoning + $36\%$ misunderstanding. OpenMedCalc~\citep{goodell2025clinical}: ``only interpretation errors were identified.'' Acurai~\citep{wood2024acurai}: $100\%$ hallucination elimination on RAGTruth ($95\%$ CI $91$--$100\%$)---strong empirical, not by-construction.

\textbf{Pattern C --- Moderate reductions ($60$--$80\%$) with residuals shifting outside the target class.} Legal RAG~\citep{dantart2026fabrication}: fabricated citations $>30\%$ $\teArrow$ $<0.2\%$. GPT-5 SimpleQA with web access~\citep{openai2025gpt5}: $47\% \teArrow 9.6\%$ (inter-condition, not within-condition). CRITIC~\citep{gou2024critic}: toxic generation $-79.2\%$.

\paragraph{The 28 GOLD citations, organised by axis and pattern.}
\begin{center}
\small
\begin{tabular}{p{2.2cm}p{4.5cm}p{3.5cm}p{1.5cm}}
\toprule
\textbf{Axis} & \textbf{Citation \& targeted class} & \textbf{Pre $\teArrow$ Post} & \textbf{Pattern}\\
\midrule
Verification & \citet{ren2025deepseekprover} --- Invalid Lean proof & any $\teArrow$ 0\% by constr. & A\\
Format & \citet{openai2024structured} --- JSON schema violation & $>$60\% $\teArrow$ 0\% by constr. & A\\
Format & \citet{zhang2023tooldec} --- Tool-call syntax & 21--100\% $\teArrow$ 0\% & A\\
Format & \citet{suresh2025dingo} --- JSON parse failure & 13--82\% $\teArrow$ 0\% & A\\
Format & \citet{dong2024xgrammar} --- Multi-format errors & 20--38\% $\teArrow$ 0\% & A\\
Format & \citet{dong2026xgrammar2} --- Malformed tool calls & 33--78\% $\teArrow$ 0\% & A\\
Arithmetic & \citet{chen2022pot} --- Calculation on GSM8K & 30\% $\teArrow$ 0\% of failures & B\\
Arithmetic & \citet{goodell2025clinical} --- Clinical arithmetic & ``only interpretation errors'' & B\\
Knowledge/RAG & \citet{wada2025radiology} --- RAG hallucinations & 8\% $\teArrow$ 0\% ($p=0.012$) & B\\
Knowledge/RAG & \citet{wood2024acurai} --- Context-conflict hallucinations & 100\% $\teArrow$ 0\% on subset & B\\
Knowledge/RAG & \citet{dantart2026fabrication} --- Fabricated legal citations & $>$30\% $\teArrow$ $<$0.2\% & C\\
Code Exec & \citet{wen2024fixing} --- \texttt{SyntaxError} (HumanEval) & 5.76\% $\teArrow$ 0.01\% & A\\
Code Exec & \citet{li2022alphacode} --- False-positive submissions & 62\% $\teArrow$ 4\% & B\\
Code Exec & \citet{shi2024code} --- 5 code-bug classes & 100\% repair (5/6 classes) & B\\
Knowledge/RAG & \citet{gao2023citations} --- Citation grounding (ELI5) & $\teApprox$50\% w/o complete support & C\\
Knowledge/RAG & \citet{openai2025gpt5} --- Factual errors & 47\% $\teArrow$ 9.6\% w/web & C\\
Knowledge/RAG & \citet{zakka2024almanac} --- Clinical citation errors & 44\% $\teArrow$ 9\% & C\\
Arithmetic & \citet{wangb2025medrac} --- Arithmetic in medical reasoning & 426 $\teArrow$ 74 errors & B\\
Code Exec & \citet{wen2024fixing} --- \texttt{NameError} repair & 22.7\% $\teArrow$ 2.3\% & C\\
Perception & \citet{gou2025uground} --- GUI grounding errors & 16.2\% $\teArrow$ 73.3\% acc. & B\\
Perception & \citet{xie2025jedi} --- GUI task failures & 5\% $\teArrow$ 27\% SR ($5.4\times$) & B\\
Perception & \citet{cheng2024seeclick} --- Element mis-location & 5.2\% $\teArrow$ 53.4\% ($10.3\times$) & B\\
Perception & \citet{liu2023deplot} --- Chart-reading errors & 38.2\% $\teArrow$ 67.6\% (+29.4 pp) & B\\
Verification & \citet{gou2024critic} --- Toxic generation & $-79.2\%$ & C\\
Verification & \citet{wangp2024mathshepherd} --- Reasoning step errors & 28.6\% $\teArrow$ 43.5\% w/rerank & B\\
Knowledge/RAG & \citet{asai2024selfrag} --- Unsupported facts & 55.5\% $\teArrow$ 22\% error & B\\
Perception & \citet{lu2024omniparser} --- Icon/element errors & 70.5\% $\teArrow$ 93.8\% (79\% red.) & B\\
Knowledge/RAG & \citet{mallen2023popqa} --- Long-tail entity errors & $\teApprox$80\% $\teArrow$ $\teApprox$50\% failure & B\\
\bottomrule
\end{tabular}
\end{center}

\paragraph{Capabilities are coarser than error classes.}
Several GOLD citations reveal ``two-for-one'' reductions where one capability addresses multiple named clusters: a single Python interpreter removes the execution-error component of arithmetic, unit conversion, simple counting, list manipulation, and date arithmetic; constrained decoding eliminates by construction both format violations and the structural component of ``missing required element''; code execution feedback strongly reduces \texttt{SyntaxError}, \texttt{NameError}, and most \texttt{TypeError} together; RAG strongly reduces factual hallucinations, fabricated citations, and outdated information jointly. The practical capability library required is therefore smaller than the count of named error categories.

\paragraph{The twelve named clusters.}
\begin{center}
\small
\begin{tabular}{p{2.5cm}p{2.2cm}p{4cm}p{2.5cm}p{1cm}}
\toprule
\textbf{Cluster} & \textbf{Axis} & \textbf{Intervention} & \textbf{Before $\teArrow$ After} & \textbf{Patt.}\\
\midrule
A. Arithmetic & Arithmetic & PAL Python~\citep{gao2022pal} & 20.1\% $\teArrow$ 61.5\% & B\\
B. Unit conversion & Arithmetic & Folded into A (Python) & --- & ---\\
C. Counting & Arithmetic & Explicit counter (gap) & --- & gap\\
D. Format/schema & Format & DINGO/SLOT~\citep{suresh2025dingo,wang2025slot} & up to $+68$ pp; 99.5\% acc. & A\\
E. Code logic/type & Code Exec & Reflexion/AgentCoder~\citep{shinn2023reflexion,huang2023agentcoder} & 80\% $\teArrow$ 96.3\% pass@1 & A/B\\
F. Multi-hop drift & Knowledge & Entity-grounded rewriter & $\teApprox$5--10 pp & C\\
G. Reasoning step & Verification & Math-Shepherd~\citep{wangp2024mathshepherd} & 28.6\% $\teArrow$ 43.5\% & B\\
H. Spec misinterpret. & Verification & Clarification~\citep{niwa2024ambignlg} & ROUGE-L $+15$ & C\\
I.a Struct.\ missing & Format & Subsumed by D (constr.\ decode) & --- & A\\
I.b Semantic missing & Verification & Subsumed by G/H & --- & C\\
J. Hallucination & Knowledge & RAG~\citep{wood2024acurai} & non-zero $\teArrow$ 0\% & B/C\\
K. Refusal & Verification & POROver~\citep{karaman2024porover} & 57.6\% $\teArrow$ 82.1\% & B\\
L. Tool/API & Format+Verif. & SAGE-Agent~\citep{suri2025clarification} & 36.5\% $\teArrow$ 65.2\% & B\\
\bottomrule
\end{tabular}
\end{center}

Eight of twelve categories have a strong citation with double-digit percentage-point improvement; three (B, C, I) lack a clean single-cluster ablation; one (F) has consistent medium-strength evidence. Category I (missing required elements) splits mechanistically into I.a (structural, absorbed by D's constrained decoders) and I.b (semantic, absorbed by G/H). Categories B (units) and C (counting) remain technical gaps: B is naturally folded into A; C is the smallest residual gap.

\paragraph{Additivity and its limits.}
\citet{patel2026sixsigma} demonstrate that stacking interventions targeting orthogonal failure modes can produce large compound reliability gains: their parallel-consensus framework yields a $14{,}700\times$ improvement over single-pass baseline---evidence for compound gains from decomposition plus consensus aggregation under their specific parallel-voting regime, not direct demonstration that heterogeneous cluster-targeted interventions compose additively without voting. \citet{shang2024agentsquare} similarly show supra-additive gains when modules target distinct failure modes. \citet{le2026interaction} reports that schema-level and prompt-level instructions interact \emph{non-additively} when sharing a prompt channel: interventions on orthogonal processing layers (decoding constraint vs.\ retrieval vs.\ training-signal vs.\ inference-time tool call) compose near-additively, while those sharing a channel may interfere.

\paragraph{The irreducible-semantic residual.}
Of the 17 ErrorAtlas categories, 13 are addressed under Patterns A, B, or C by one of the six capability axes; four are not---residuals of inappropriate refusal beyond preference optimisation, specification misinterpretation not closed by clarification, problem-decomposition reasoning bottlenecks, and a ``user wanted something different'' semantic remainder. These are classes where the failure is in choosing what to do, not executing it. Proposition~1's prediction of $O(\log)$ rather than $O(0)$ residual reliability reflects exactly this irreducible core. The framework does not claim $100\%$ coverage of all failures; it claims polylog-bounded \emph{capability-eliminable} failures, with the named semantic residual as the floor.

\section{Counter-Evidence Re-Audits}
\label{app:reaudits}

Five prominent papers are routinely cited as evidence that LLM reliability decays steeply with length~\citep{dziri2023faith,kuratov2024babilong,kwa2025horizon,wan2026fano,karpinska2024nocha}. A careful re-reading shows that \emph{every one} of these papers decays over a variable distinct from raw token length $n$. Identifying the decay axis is not the same as dissolving the practical concern: where compositional graph size, fact count, or evidence scope grow with problem length, the framework predicts steep failure curves. The contribution is to identify which interventions help (capability provisioning along the actual decay axis) and which do not.

\paragraph{\citet{dziri2023faith}, Faith and Fate.}
GPT-4 multi-digit multiplication accuracy drops from 59\% (3-digit) to 4\% (4-digit) zero-shot, with the authors theorising ``probability of incorrect predictions converges exponentially to $\teApprox 1$ for abstract compositional tasks.'' The decay variable is \emph{compositional graph size $N$}, not raw token length: the $3 \times 3$ graph has on the order of $d^2$ partial products plus carries, and multi-digit multiplication is engineered so that $k_{\text{hard}} \teApprox N$---every node in the computation graph is a hard decision. For natural-language tasks where $k_{\text{hard}} \ll n$, Dziri's regime is the \emph{boundary case} in which our framework reduces to their result. A clean $(1-\varepsilon)^N$ exponential cannot simultaneously reproduce 59\% at 3-digit and 4\% at 4-digit for any single per-node $\varepsilon$---the observed drop is locally steeper than per-node iid exponential, consistent with a finite catalogue of failure modes exhausting as $N$ grows. \emph{Where the framework concedes}: in adversarial compositional tasks engineered so $k_{\text{hard}} \teApprox N$, the framework reduces to Dziri's regime and does not relieve it; the relocation is informational, not magical.

\paragraph{\citet{kuratov2024babilong}.}
The abstract reports ``performance declines sharply with increased reasoning complexity'' and models ``effectively utilise only $10$--$20\%$ of the context.'' Read in detail, BABILong varies two axes: context length $n$ ($0$K to $10$M tokens) and number of supporting facts $k$ (QA1 $=1$ fact to QA3 $=3$ facts). The sharp decline is in $k$, not $n$. For QA1 (single-fact), most models ``perform well up to 4,000 tokens''---a plateau, not exponential decay. The famous RAG-flat-across-length result ($60\%$ on single-fact QA \emph{independent of context length}) is the cleanest possible demonstration that when relevant evidence is in window, length does not matter. Recurrent Memory Transformers maintaining performance to 50M tokens further confirms effective $k$ is determined by architecture, not raw $n$. \emph{Where the framework concedes}: for multi-hop tasks whose required fact count grows with task complexity, BABILong's sharp $k$-axis decay is the framework's \emph{predicted} trouble, not a counter-example.

\paragraph{\citet{kwa2025horizon} (METR).}
Per-model success fits a logistic in $\log(\text{human-task-duration})$: $S(t) = \sigma(\beta (\log t - \log h))$. Logistic-in-log-length is \emph{mathematically sublinear} in length itself. The 80\%-horizon being $4$--$6\times$ shorter than the 50\%-horizon is a steep within-model cliff but compatible by construction with the polylog result---a cliff at a specific capacity threshold is a manifold-transition signature. The famous exponential---capability doubling every seven months---is an \emph{inter-model} claim about how the horizon $h$ moves across generations, orthogonal to within-model decay shape. \emph{Where the framework concedes}: the within-model logistic-in-log cliff is steep on a practical scale; calling it ``sublinear in length'' is technically correct but engineering-useful only for systems designed to operate well below the 50\% horizon.

\paragraph{\citet{wan2026fano}, Fano-style upper bound.}
This paper theorises super-linear information-demand growth and identifies an ``accuracy cliff'' at capacity overflow: ``when the task's information demand surpasses the model's output capacity, performance does not degrade gracefully but instead collapses sharply.'' The behavioural prediction is a \emph{threshold}, not smooth $(1-\varepsilon)^n$. A cliff at a specific capacity threshold is exactly the manifold-transition behaviour our framework predicts at the boundary between covered and uncovered modes; Wan et al.'s mechanism (capacity overflow) is one specific cause, complementary to ours. \emph{Where the framework concedes}: Wan documents a mechanism by which $|C|$ effectively exceeds the model's coverage in a single forward pass; Pattern A interventions (constrained decoding, formal verification) are the kind of response the framework expects but does not automatically supply.

\paragraph{\citet{karpinska2024nocha}.}
GPT-4o achieves 55.8\% pair accuracy across 1,001 minimally-different true/false claim pairs about long fictional books (mean length 127K tokens). The paper reports performance by \emph{evidence scope}, not by context length: 59.8\% on sentence-level retrieval, 47.6\% on passage-level, 41.6\% on global reasoning. No per-context-length curve is reported. The decay axis is the number of evidence pieces that must be integrated---a $k$-axis quantity, not an $n$-axis one. NoCha is therefore not counter-evidence to a sublinear-in-$n$ claim. \emph{Where the framework concedes}: evidence-scope decay is a $k$-axis observation that the framework predicts cannot be addressed by scaling raw context length; it requires retrieval or process supervision along the actual decay axis.

\paragraph{Unifying observation.}
Every counter-paper decays over a variable that is not $n$: compositional graph size, fact count, log-time horizon, capacity threshold, or evidence scope. The apparent rapid decay is, in each case, in a quantity our framework already concentrates the action in ($k_{\text{hard}}$ and $|C|$), not in raw sequence length. This is a relocation, not a dissolution. Where $k_{\text{hard}}$ grows with task length---adversarial compositional structure, multi-hop fact chains, long horizons that force more decisions---reliability remains hard; the framework's value is directing intervention toward capability provisioning along the actual decay axis rather than toward context-window or compute-budget expansion that does not help.

\newpage
\input{checklist.tex}

\end{document}
